% =====================================================================
%  SVD-TA / Compressed Traffic Assignment -- manuscript v3
%  Numerical-rerun pass: 2026-07-24   (repro code: CompressedTAP @ 2adc6d4, branch version_3)
%  Re-evaluated this pass : Table 2 (Sioux consistency); Table 4 Panel A -- Chicago Sketch
%                           (all tau) and Philadelphia (tau = 0, 0.67).
%  NOT re-evaluated (--)  : Panel A Chicago Regional; Table 5 (rank); Table 4 Panel B.
%  C++/Python metrics verified CONSISTENT on the Sketch instance
%  (Gap_F within 0.03pp, R^2 within 1.2e-3).
% =====================================================================
\documentclass[11pt]{article}
\usepackage{amsfonts,amsmath,amsthm,amssymb}
\usepackage{graphicx}
\usepackage{booktabs}
\usepackage{float}
\usepackage{longtable}
\usepackage{geometry}
\usepackage[table]{xcolor}
\usepackage{url}
\usepackage[normalem]{ulem}

% Conservative revision markup: unchanged submitted text remains black;
% proposed replacements/additions are red.
\newcommand{\revnote}[1]{\par\noindent{\color{red}\textbf{Revision note:} #1}\par}


\newtheorem{exmp}{Example}[section]
\newtheorem{theorem}{Theorem}[section]
\newtheorem{lemma}[theorem]{Lemma}
\newtheorem{proposition}[theorem]{Proposition}
\newtheorem{remark}{Remark}[section]
\newtheorem{corollary}[theorem]{Corollary}
\newtheorem{definition}[theorem]{Definition}
\newtheorem{conjecture}[theorem]{Conjecture}
\newtheorem{axiom}[theorem]{Axiom}
\newtheorem{assumption}{Assumption}[section]

\title{Compressed Traffic Assignment with the \\Augmented Lagrangian Method\footnote{
This work was carried out at the Fulton School of Engineering, Arizona State University, Tempe, AZ. X. Zhou and Y. Li are with School of Sustainable Engineering and the Built Environment, ASU. P. Li is with Norfolk Southern Corporation. D. Bertsekas is with School of Computing and Augmented Intelligence, ASU.}}
\author{Xuesong Zhou, Peiheng Li, Yuchao Li, and Dimitri Bertsekas}
\date{April 2026}

% ---- draft version stamp (bump DOCVER on every circulated build) ----
\newcommand{\docver}{v3.6}
\newcommand{\docdate}{2026-07-25}
\newcommand{\doccommit}{0077beb}
\newcommand{\docstamp}{Draft \docver\ \textbullet\ built \docdate\ \textbullet\ code
  \texttt{CompressedTAP@\doccommit} (\texttt{version\_3}) \textbullet\ numerical section
  rewritten; flow-weighted basis; paired-comparison figures; body in black, red = notes to co-authors}

% Author-facing notes (red). \shownotesfalse suppresses them for a submission build;
% the calls can then stay in the source.
\newif\ifshownotes \shownotesfalse
\newcommand{\conote}[1]{\ifshownotes{\color{red}\small\textsf{[Note to co-authors: #1]}}\fi}
\newcommand{\conoteY}[1]{\ifshownotes{\color{red}\small\textsf{[Note to Yuchao: #1]}}\fi}

\begin{document}
{\Large\bfseries Compressed Traffic Assignment --- Numerical Section}\par
\vspace{3pt}
{\small\textsf{\docstamp}}\par
\vspace{3pt}
{\small Standalone extract of the rewritten computational study. Body text is black;
\textcolor{red}{red marks notes to co-authors} only, to be removed before submission.
Every number traces to a CSV in \texttt{repro\_v3/results}.}
\vspace{8pt}\hrule\vspace{8pt}
\section{Computational Studies}
\label{sec:experiments}
\label{sec:experiments}\label{sec:computational}


The computational study answers one question: \emph{when does compressing the minor-path
variables actually reduce solution time, and by how much?} The submitted version reported
threshold and rank sweeps on three road networks. Re-running those experiments exposed a
sharper structure than the sweeps alone revealed, and the study is organized around it here:
a cost model that predicts when compression pays (\S\ref{sec:cost-model}); a controlled grid
family that isolates the two variables the model identifies (\S\ref{sec:grid-study}); the
rank, which turns out to be a speed parameter rather than an accuracy parameter
(\S\ref{sec:rank}); the criterion selecting the retained subspace, where weighting by the
nominal flow halves the objective gap on route-rich instances at no additional cost
(\S\ref{sec:weighted-basis}); the affine offset $w_0$, without which the representation does not work
(\S\ref{sec:offset-study}); the solution operator, where the benefit turns out to be specific
to the augmented Lagrangian (\S\ref{sec:operator-study}); and the road networks, where the
model's predictions are tested against practice (\S\ref{sec:road-networks}).
Section~\ref{sec:repeated} then runs the repeated-scenario setting the framework is actually
aimed at, building the basis once and reusing it across perturbed scenarios, against a
warm-started full-path solve, and reports how many related solves repay the preprocessing.
Section~\ref{sec:scope-limits} states the scope and the negative results.

\conote{This section is a full rewrite rather than an update of the submitted one, so it is
set in black and red is reserved for these notes. The organizing change is that it leads with
a cost model and then tests it, instead of reporting parameter sweeps. Referee~1's own
diagnosis --- that the dense $(n-s)\times r$ nonnegativity block caps the gain --- is exactly
what the model formalizes, which is deliberate.}


\subsection{Experimental Setup}
\label{sec:setup}


\textbf{Instances.} Two families are used. The \emph{grid family} is an $N\times N$
Manhattan grid with bidirectional links, unit free-flow time, capacity $600$, and BPR
parameters $(\alpha,\beta)=(0.15,4)$; demand runs from the left column to the right column
with $720/N$ vehicles per OD pair, so the mid-cut volume-capacity ratio stays near $1.2$ at
every $N$ and congestion does not drift as the grid grows. Candidate paths are generated by
a deterministic penalty $K$-shortest-path routine at depth $K$. The \emph{road networks} are
the three instances of the submitted study, Chicago Sketch, Chicago Regional, and
Philadelphia (Table~\ref{tab:network-characteristics}), together with Sioux Falls and a
sequence of progressively enriched Chicago Sketch pools used to vary route richness at fixed
network and demand.

\textbf{Bases.} Every compressed instance is solved with both subspace criteria --- the plain
singular subspace of $B_2$ and the flow-weighted subspace of
$\operatorname{diag}(\sqrt{w_0})B_2$ (Section~\ref{sec:weighted-basis}) --- run back to back, so
each comparison is within-session and unaffected by machine load.

\textbf{Implementations.} All experiments are run in two independent implementations of the
same algorithm: a Python/NumPy reference and a compiled C\texttt{++} solver. Reported
speedups are always \emph{within-implementation} ratios --- an implementation's compressed
time divided into its own uncompressed time on the same instance --- because, as
\S\ref{sec:scope-limits} documents, the ratio depends on the relative throughput of dense and
sparse kernels and is not portable between implementations.

\textbf{Protocol.} Single thread, matched stopping tolerance $10^{-4}$ on a common relative
duality-gap certificate, quiet machine. Every terminal iterate is converted to a feasible
point by the same per-OD projection before any objective is compared; $\delta_F$ reports the
$\ell_1$ mass that conversion moves and $\mathrm{Gap}_F$ the resulting feasible objective
gap against the best feasible point found by any method on that instance. Timings on the
smaller instances are medians of three runs; on the largest instances they are single runs,
and repeated identical solves varied by as much as $40\%$ under machine load. It is worth
stating this precision explicitly, because it governs which comparisons the study is entitled
to make. Re-measuring identical solves --- same instance, same settings, same binary --- across
the campaign gave spreads of $1.1$ to $1.6\times$ typically and $2.8\times$ in the worst case,
so we treat $1.45\times$ as a conservative noise floor and read no speed difference smaller
than that as an effect. Feasible objective gaps carry no such caveat: they are deterministic
and reproduced bit for bit across every repetition, so accuracy comparisons are exact even when
the corresponding timings are not.

\conote{Two accounting decisions are ours to make before submission. (i) None of the reported
speedups charge the SVD or the nominal-flow computation; Referee~1 objected to exactly this
and the objection is currently unanswered. Proposal: report every ratio twice, online-only and
fully charged, with a break-even count in related solves. (ii) Large-instance timings are
single runs at $\pm40\%$; medians of three would cost roughly two machine-hours and would let
the near-parity entries be read as ordered.}


\subsection{When Compression Pays: A Cost Model}
\label{sec:cost-model}


Compression replaces $n$ path variables by $s+r$, but the gradient of the compressed problem
is not simply $(s+r)/n$ times cheaper. Each evaluation must reconstruct the minor flows to
enforce $w_0+U_rz\ge 0$, and $U_r$ is dense: the reconstruction and the corresponding
gradient term each cost about $(n-s)\,r$ operations, so the compressed gradient pays
approximately $2r$ dense operations \emph{per minor path}. The uncompressed gradient, by
contrast, pays the sparsity of a path --- the number of links it contains --- per path. The
comparison is therefore not between $n$ and $s+r$ but between
\begin{equation}
\label{eq:cost-rule}
2r \quad\text{and}\quad \text{(average links per path)} .
\end{equation}
Compression reduces work per evaluation when $2r$ is at most the average path length, and
adds work when it is not. Two consequences are tested in what follows. First, the benefit
should grow with the size of the candidate pool at fixed network, because the uncompressed
side grows with the pool while the compressed side does not. Second, the rank should behave
mainly as a cost parameter, since it enters \eqref{eq:cost-rule} linearly but affects the
represented subspace only weakly once the major paths carry the bulk of the flow.


\subsection{A Controlled Study on Grid Networks}
\label{sec:grid-study}


The grid family isolates the two quantities in \eqref{eq:cost-rule}: the grid dimension $N$
sets the average path length ($\approx N$ links), and the generation depth $K$ sets the pool
size. Table~\ref{tab:grid-matrix} reports the speedup of the compressed augmented Lagrangian
over the uncompressed one at fixed threshold ($75$th percentile of the nominal flow, giving
$\approx 74\%$ variable reduction) and rank $r=20$, in the C\texttt{++} implementation.

\begin{table}[h]
\centering
\caption{Speedup of compressed over uncompressed ALM on the grid family
(C\texttt{++}, $r=20$, $\approx74\%$ variable reduction, matched tolerance). Speedup rises
with pool depth $K$ at every grid size. Feasible objective gaps are at most $0.55\%$ in the
$K\ge32$ columns and $1.4$--$3.4\%$ in the $K=4$ column, where the fixed threshold removes
genuine flow from a thin pool.}
\label{tab:grid-matrix}
{\small
\begin{tabular}{lrrrrr}
\toprule
& $K=4$ & $K=8$ & $K=16$ & $K=32$ & $K=64$\\
\midrule
$N=8$  & 0.55 & 0.59 & 1.53 & 1.45 & 2.01\\
$N=16$ & 2.09 & 0.89 & 1.75 & 5.16 & \textbf{6.82}\\
$N=24$ & 2.25 & 1.18 & 0.73 & 3.87 & 4.86\\
$N=32$ & 1.76 & 2.76 & 0.96 & 3.13 & 5.02\\
\bottomrule
\end{tabular}}
\end{table}

The prediction holds (Figure~
ef{fig:speedup-mechanisms}a). Reading across any row, the speedup rises with pool depth, reaching
$6.82\times$ at $N=16$, $K=64$ with a feasible objective gap of $0.0000\%$; the $K=32$ and
$K=64$ columns are uniformly between $3.1\times$ and $6.8\times$ with gaps at most $0.55\%$.
Reading down the columns, the effect strengthens with grid size once the pool is deep enough,
consistent with longer paths making the right-hand side of \eqref{eq:cost-rule} larger. The
$K=8$--$16$ interior is noisy because those cells solve in under ten seconds. The $K=4$
column carries $1.4$--$3.4\%$ gaps: with only four candidate paths per OD pair, a fixed
$75$th-percentile threshold classifies genuinely used paths as minor.

\conote{Open question for the group: the grid is where the largest speedups live, but it is
synthetic, and our own verification found the C\texttt{++} uncompressed baseline on grids
converging to an objective $4.6$--$4.9\%$ above the Python one. Until that baseline is
tightened, a referee can argue these ratios compare two under-converged solves. Do we tighten
the baseline, or present the grid strictly as a mechanism study and let the road networks
carry the efficiency claim?}



\begin{figure}[h]
\centering
\includegraphics[width=\textwidth]{fig/fig_speedup_mechanisms.pdf}
\caption{The two mechanisms that govern the compression speedup. (a) On the grid
family, the gain grows with the depth of the candidate pool at every grid size, because the
uncompressed problem grows with the pool while the compressed dimension does not; the dashed
line marks parity. (b) On Sioux Falls the rank moves the speedup by close to an order of
magnitude in both implementations while the feasible objective gap, annotated at each point,
is unchanged to four decimal places --- over this range the rank is a speed parameter, not an
accuracy parameter.}
\label{fig:speedup-mechanisms}
\end{figure}

\subsection{The Rank is a Speed Parameter}
\label{sec:rank}


Equation~\eqref{eq:cost-rule} predicts that the rank controls cost linearly. The submitted
study swept $r\in\{50,100,150,200\}$ and found solution quality insensitive to $r$ while cost
grew --- correct, but sampled entirely on the expensive side. Table~\ref{tab:impact-rank}
sweeps downwards on Sioux Falls in both implementations.

\begin{table}[h]
\centering
\caption{Effect of the rank on Sioux Falls ($\tau=600$, $91\%$ variable
reduction, matched tolerance, medians of three runs). Speedup is against each
implementation's own uncompressed baseline ($1.34$s Python, $8.20$s C\texttt{++}). The
feasible objective gap is unchanged to four decimal places across the entire range.}
\label{tab:impact-rank}
{\small
\begin{tabular}{llrrr}
\toprule
Implementation & $r$ & CPU (s) & \textbf{Speedup} ($\times$) & Feasible obj.\ gap (\%)\\
\midrule
Python & 10 & 0.13 & \textbf{10.31} & 3.4388\\
       & 20 & 0.37 & \textbf{3.62}  & 3.4389\\
       & 50 & 1.15 & \textbf{1.17}  & 3.4386\\
\midrule
C\texttt{++} & 10 & 0.97 & \textbf{8.45} & 0.4541\\
             & 20 & 1.20 & \textbf{6.83} & 0.4541\\
             & 50 & 2.88 & \textbf{2.85} & 0.4541\\
\bottomrule
\end{tabular}}
\end{table}

The speedup varies by nearly an order of magnitude while the feasible objective gap moves in
the fourth decimal place (Figure~
ef{fig:speedup-mechanisms}b). Over this range the rank is therefore almost purely a speed
parameter: it multiplies the dense reconstruction without moving the solution. The optimum is
interior rather than at $r\to0$, because too small a subspace eventually lengthens the inner
solves --- visible on the grid, where $N=16$, $K=32$ gives $5.24\times$, $7.34\times$ and
$5.44\times$ at $r=10$, $20$ and $40$. Combining this with \eqref{eq:cost-rule} gives the
design rule used throughout the rest of the study: \emph{choose $2r$ no larger than the
average path length}, which on these instances means $r\approx10$--$20$ rather than the
$r=50$ of the submitted experiments. Section~\ref{sec:road-networks} quantifies what that
choice is worth on a road network.

\conoteY{This rule is the direct answer to Referee~2's request for analytical guidance on
$r$, and it is currently the only place we give a parameter rule. He asked the same question
about $\tau$ --- uniform or per-OD, and how it relates to convergence --- and we still choose
$\tau$ by quantile with no argument. If you can supply even a heuristic derivation for $\tau$
analogous to \eqref{eq:cost-rule}, it would close his second objection as cleanly as this
closes the first.}


\subsection{Which Subspace to Keep: Flow-Weighted Selection}
\label{sec:weighted-basis}

Section~\ref{sec:rank} fixes how large the retained subspace should be. A separate question is
which subspace of that size to retain. The natural answer, and the one used above, is the
leading left-singular subspace of $B_2$, which minimises $\|B_2 - \Pi B_2\|_F$. That criterion
counts every minor path equally, and it is not the criterion the problem asks for: a minor path
enters the link load in proportion to the flow it carries, so an error on a path carrying a
thousand vehicles matters a thousand times more than the same error on a path carrying one.

We therefore scale the rows of $B_2$ by the square root of the nominal flow before
factorising. Precisely, with $\mathcal M$ the minor set, $n_m=|\mathcal M|$, and
$B_2\in\mathbb R^{n_m\times m}$ the minor path-link matrix, set
\begin{equation}
\label{eq:weighted-basis}
S=\operatorname{diag}(s),\qquad s_p=\sqrt{w_{0,p}}+\varepsilon,\qquad
\Phi_r=\widetilde U_r,
\end{equation}
where $\widetilde U_r\in\mathbb R^{n_m\times r}$ collects the $r$ leading left singular
vectors of $S\,B_2$ ordered by singular value, and the minor flows are reconstructed as
$x_{\mathcal M}=w_0+\Phi_r z$ exactly as before. The constant $\varepsilon=10^{-9}$ keeps $S$
invertible when a nominal flow vanishes, which is not a rare case: on Sioux Falls $80.6\%$ of
minor paths carry no nominal flow, while on the Chicago Sketch pools every minor path carries
some, so the two instance families sit in genuinely different regimes.

It is worth being exact about what \eqref{eq:weighted-basis} does and does not achieve, because
the natural reading is stronger than the truth. Row scaling makes high-flow paths dominate the
factorisation, which is the intent. It does \emph{not} return the subspace minimising the
flow-weighted reconstruction loss $\sum_p w_{0,p}\|(B_2-\Pi B_2)_p\|^2$: by Eckart--Young that
minimiser is $\operatorname{span}(S^{-1}\widetilde U_r)$, obtained by mapping the singular
vectors back out of the scaled coordinates, whereas \eqref{eq:weighted-basis} uses them
directly as path-flow coordinates. The gap is measurable. Evaluating
$\min_Z\|S(B_2-\Phi Z)\|_F$ at $r=20$ gives $249.1$ for \eqref{eq:weighted-basis} against
$209.1$ for $S^{-1}\widetilde U_r$ and $227.9$ for the plain singular subspace on Sioux Falls,
and $404.3$, $393.7$ and $394.7$ respectively on Chicago Sketch; the two subspaces are nearly
orthogonal on Sioux Falls, where most nominal flows vanish, and close on Chicago Sketch, where
none do. So \eqref{eq:weighted-basis} is a row-scaled heuristic, justified below by what it
achieves on the quantity that matters --- the feasible objective gap of the solved problem ---
rather than by optimality in the reconstruction norm. Characterising the basis that is optimal
for the solved problem, as opposed to for reconstruction, is open.

The two factorisations are the same size and cost the same --- $0.94$s against $0.97$s on
Chicago Sketch at $r=20$ --- so any difference in solution time is a convergence effect and not
preprocessing overhead. Table~\ref{tab:weighted-basis} varies the candidate pool at fixed
network, demand, threshold and rank.

\begin{table}[h]
\centering
\caption{Effect of flow-weighted subspace selection on Chicago Sketch ($\tau=4.54$, $r=20$,
matched tolerance). Network, demand and parameters are fixed; only the richness of the
candidate pool varies. Both bases are solved back to back on each instance.}
\label{tab:weighted-basis}
{\small
\begin{tabular}{rrrrrr}
\toprule
& \multicolumn{2}{c}{flow-weighted} & \multicolumn{2}{c}{plain} & \\
\cmidrule(lr){2-3}\cmidrule(lr){4-5}
$\bar K$ & speedup & gap (\%) & speedup & gap (\%) & gap reduction\\
\midrule
2.45  & \textbf{1.62} & \textbf{0.0284} & 1.55 & 0.0291 & $-2.4\%$\\
3.49  & \textbf{3.71} & \textbf{0.0202} & 3.34 & 0.0232 & $-12.9\%$\\
8.64  & \textbf{1.55} & \textbf{0.2086} & 1.34 & 0.7024 & $\mathbf{-70.3\%}$\\
11.09 & \textbf{0.83} & \textbf{0.4350} & 0.66 & 1.1914 & $\mathbf{-63.5\%}$\\
15.34 & \textbf{1.20} & \textbf{0.8649} & 1.15 & 1.8651 & $\mathbf{-53.6\%}$\\
\bottomrule
\end{tabular}}
\end{table}

\textbf{The effect is on accuracy, and only on accuracy.} The gap column improves at every
richness point, by more than half wherever $\bar K \ge 8.64$, and these figures are exact:
feasible objective gaps are deterministic and reproduce bit for bit across runs. The speedup
column moves in the same direction, but by margins of $4$ to $26\%$, which is inside the
measurement noise reported in Section~\ref{sec:setup}. Pooling all $34$ controlled
weighted-versus-plain pairs measured in this study --- each a single instance with both bases
solved back to back --- the gap improves in $27$, whereas the speedup improves by more than the
noise floor in only $2$. We therefore claim an accuracy improvement and explicitly do not claim
a speed improvement.

The trend across the table is the explanation: a rich candidate pool is rich precisely in paths
that carry almost no flow, and the unweighted criterion spends its rank budget representing
them. This is why the regime that was least favourable in Section~\ref{sec:road-networks} ---
rich pools, where the compressed problem showed the largest gap --- is the regime the weighting
repairs most.

The grid family reproduces the effect on the axis that can be resolved. Across its twenty cells
the weighted basis lowers the feasible gap in $15$, taking the mean from $1.285\%$ to
$1.112\%$, with the largest improvements again at high pool depth (at $N=32$, $K=64$ the gap
falls from $0.406\%$ to $0.000\%$). Sioux Falls, at $91\%$ variable reduction, shows the same
sign and the largest relative movement, $11$ to $20\%$. The conclusion is that weighting makes
the subspace criterion consistent with the objective --- the retained directions are the ones
that carry flow --- and that this buys accuracy at no preprocessing cost, not speed.

\subsection{The Affine Offset is Load-Bearing}
\label{sec:offset-study}


The representation is affine, $x_{\text{minor}}=w_0+U_rz$, and the offset $w_0$ is the
nominal minor flow. Because $U_r$ is signed, the feasible set $\{w_0+U_rz\ge0\}$ contains the
starting point $z=0$ in its interior only by virtue of $w_0$; without the offset, $z=0$ lies
on the boundary of every reconstructed row simultaneously. To test whether this matters we
remove the offset exactly --- setting $w_0=0$ while moving its flow into the base link volume
and the effective demand, so the model is unchanged --- and re-solve.

The effect is large. Across the twenty grid cells of Table~\ref{tab:grid-matrix} the mean
speedup falls from $2.60\times$ to $1.07\times$ and the mean feasible gap grows from $0.75\%$
to $2.42\%$; the cells where compression is strongest collapse ($N=24$, $K=64$ from
$4.86\times$ to $0.45\times$; $N=16$, $K=64$ from $6.82\times$ to $1.86\times$). On Sioux
Falls the compressed solve slows from $1.7$s to $7.9$s, a fall from $6.2\times$ to
$1.34\times$, with the cost per inner iteration rising fourfold at nearly unchanged iteration
count: with the hinge active at the iterate, the augmented Lagrangian subproblem becomes
nonsmooth exactly where the solver works. In the Python implementation the offset-free
variant is not slower, and the accuracy penalty appears only on the larger instances, so the
speed half of this conclusion is specific to the compiled solver while the accuracy half is
not. The offset is therefore not an implementation convenience: \emph{a signed subspace of
nonnegative flows is usable only relative to an interior anchor}, and $w_0$ supplies it.

The two representational choices are independent. Removing the offset from the
\emph{flow-weighted} representation of Section~\ref{sec:weighted-basis} degrades it by the same
mechanism and by a similar amount: across the grid family the mean feasible gap rises from
$1.112\%$ to $1.394\%$, and individual cells collapse in the same way (at $N=32$, $K=64$,
$5.41\times$ at a gap of $0.000\%$ becomes $1.00\times$ at $1.681\%$). Choosing the subspace
better does not remove the need for an anchor inside it, and choosing the anchor well does not
substitute for aligning the subspace with the flow. The two act on different parts of the
representation and their benefits do not overlap.

\conote{This result cuts both ways and we should decide how to present it. It explains why the
representation works, but it also quantifies the dependence Referee~1 objected to: $w_0$ is
obtained by solving the full problem, and we have now proved the method needs it. The
strongest available answer is an experiment we have not run --- reuse a \emph{stale} $w_0$
from a perturbed instance and show the speedup survives. That single test would answer the
circularity objection, the preprocessing objection and the ``which application setting''
objection at once. Recommend we run it before resubmission.}


\subsection{Which Operator Benefits from Compression}
\label{sec:operator-study}


The compressed representation is independent of the solution method, so one may ask whether
the benefit is a property of the representation or of the augmented Lagrangian. We solve the
same compressed instances with three operators: the augmented Lagrangian; gradient projection
on the compressed set; and a reduced-gradient method that eliminates OD conservation by
anchor substitution. Table~\ref{tab:operator-compression} reports the result on grid cells in
the regime where compression is effective.

\begin{table}[h]
\centering
\caption{Speedup from compression, by solution operator, on grid cells
(Python, $w_0$ representation, matched tolerance). Each column is the operator's own
uncompressed time divided by its compressed time. Frank--Wolfe has no compressed counterpart
because the linear oracle over the compressed set is itself a linear program.}
\label{tab:operator-compression}
{\small
\begin{tabular}{lrrr}
\toprule
Grid cell & ALM & Gradient projection & Reduced gradient\\
\midrule
$N=16$, $K=16$ & \textbf{1.11} & 0.00 & 0.07\\
$N=16$, $K=32$ & \textbf{1.22} & 0.01 & 0.26\\
$N=16$, $K=64$ & \textbf{1.96} & 0.01 & 0.13\\
$N=32$, $K=16$ & \textbf{3.85} & 0.02 & 0.34\\
$N=32$, $K=32$ & \textbf{1.74} & 0.04 & 0.22\\
$N=32$, $K=64$ & \textbf{2.54} & 0.10 & 0.26\\
\bottomrule
\end{tabular}}
\end{table}

Only the augmented Lagrangian benefits. The reason is structural rather than incidental. An
iteration costs gradient work over the \emph{variables} plus feasibility work over the
\emph{constraints}; compression reduces the former and, because the reconstruction is dense,
increases the latter. The augmented Lagrangian carries feasibility in the objective through
multipliers and a penalty, so its iteration cost is gradient-dominated and the variable
reduction passes through to wall-clock time. Gradient projection instead derives its speed
from a projection that is separable across OD pairs and available in closed form; on the
compressed set $\{A_1y+Mz=d_{\mathrm{eff}},\,y\ge0,\,w_0+U_rz\ge0\}$ the latent variables
$z$ appear in every OD equality and in every reconstructed bound, so separability is lost and
the projection degrades into a coupled quadratic program. On the road networks this is not
merely slower but infeasible: the exact projection requires a dense factorization of order
the number of OD pairs, which at $\ell=17{,}464$ does not scale. The reduced-gradient method
removes the OD equalities by substitution, but the binding cost is the dense reconstruction,
which substitution does not touch; what elimination removes, the representation restores.

We therefore state the scope explicitly: \emph{variable reduction accelerates a method whose
per-iteration bottleneck lies in the variables, not in the constraints.} This is the
measured justification for pairing the compressed representation with an augmented Lagrangian
rather than a projection-type method, and it is the reason the framework of
Section~\ref{sec:compression_framework} is built that way.

\conote{Referee~2 wrote that ``the augmented Lagrangian method is quite standard'' and wanted
a tailored algorithm. This subsection is our answer: not a new algorithm, but a demonstration
that the pairing is forced. Whether that satisfies him is the main uncertainty in the
revision. He also asked for comparison against state-of-the-art path-based methods; we compare
four operators but no bush-based method (Algorithm~B, OBA, iTAPAS). Do we add one, or argue
the comparison is out of scope?}


\subsection{Road Networks}
\label{sec:road-networks}


Three questions remain for the road networks: do the operators agree on the same answer; does
the richness prediction of \S\ref{sec:cost-model} hold on a real network; and what does the
rank rule of \S\ref{sec:rank} buy in practice.


\subsubsection*{Consistency of the operators}


Before comparing efficiency we verify that the methods share a numerical target.
Table~\ref{tab:operator-consistency} solves one instance with all operators under matched
tolerances and a common feasibility conversion.


\begin{table}[h]
\centering

\caption{Numerical consistency on Sioux Falls under matched stopping tolerances. The
reference is the best feasible point over the uncompressed operators (full-path gradient
projection, feasible objective $4.231\times10^{6}$, certificate $9.1\times10^{-10}$). The
augmented Lagrangian's feasibility-based stopping rule leaves it $0.048\%$ above that
reference; the near-agreement of all operators is the consistency this table quantifies.}
\label{tab:operator-consistency}
\resizebox{\textwidth}{!}{%
\begin{tabular}{llrrrrr}
\toprule
Representation & Method & Raw obj. diff. (\%) & $\|A\tilde x-d\|_\infty$ & $\delta_F$ (\%) & Feasible obj. gap (\%) & Link-flow diff. (\%)\\
\midrule
Full ($\tau=0$) & ALM & 0.048 & $8.7\times10^{-11}$ & $1.3\times10^{-9}$ & 0.048 & 1.156\\
Full ($\tau=0$) & Frank--Wolfe & $1.2\times10^{-4}$ & $2.4\times10^{-10}$ & $2.7\times10^{-12}$ & $1.2\times10^{-4}$ & $5.1\times10^{-4}$\\
Full ($\tau=0$) & Gradient projection & $9.8\times10^{-8}$ & $5.5\times10^{-11}$ & $1.5\times10^{-11}$ & $9.8\times10^{-8}$ & 0.005\\
\midrule
Compressed ($\tau=600$, $r=20$) & ALM & 3.489 & $1.7\times10^{-10}$ & $1.5\times10^{-5}$ & 3.489 & 4.679\\
Compressed ($\tau=600$, $r=20$) & Gradient projection (signed) & 3.308 & $7.8\times10^{-11}$ & 0.102 & 3.461 & 4.655\\
\bottomrule
\end{tabular}}

\end{table}


The compressed rows agree with each other to within $0.03$ percentage points of feasible
objective gap, confirming that the gap is a property of the representation at this
$(\tau,r)$ and not of the operator, while the uncompressed rows bracket the reference to
within $0.05\%$. Frank--Wolfe is the most accurate uncompressed operator at loose tolerance
and is used as an honest baseline where relevant.


\begin{table}[h]
\centering
\caption{Network characteristics. Singleton paths correspond to the OD pairs with a unique feasible path and are excluded from the compressed problem.}
\label{tab:network-characteristics}
{\small
\begin{tabular}{lrrr}
\toprule
 & Chicago Sketch & Chicago Regional & Philadelphia \\
\midrule
Nodes / Links & 933 / 2{,}950 & 12{,}982 / 39{,}018 & 13{,}389 / 40{,}003 \\
Multi-path variables ($n$) & 42{,}774 & 879{,}625 & 1{,}846{,}068 \\
Multi-path OD pairs ($\ell$) & 17{,}464 & 353{,}582 & 601{,}644 \\
Average paths per multi-path OD ($\bar K=n/\ell$) & 2.45 & 2.49 & 3.07 \\
Multi-path flow share & 23.5\% & 59.2\% & 70.9\% \\
\bottomrule
\end{tabular}
}
\end{table}

\subsubsection*{Route richness and the rank rule}


Table~\ref{tab:cutoff-summary} varies route richness on Chicago Sketch at fixed network and
demand, by enriching the candidate pool from the submitted $\bar K=2.45$ to $\bar K=15.34$,
and reports the same axis at the rank chosen by \S\ref{sec:rank} and at the submitted rank.

\begin{table}[h]
\centering
\caption{Route-richness axis on Chicago Sketch: same network and demand,
candidate pools of increasing average richness $\bar K$, $\tau=4.54$, Python implementation.
The $r=20$ column applies the rule of \S\ref{sec:rank}; the $r=50$ column is the submitted
rank, shown for comparison at identical accuracy.}
\label{tab:cutoff-summary}
{\small
\begin{tabular}{rrrrrrr}
\toprule
$\bar K$ & $\bar K/q$ & Full dim.\ $n$ & Compressed dim.\ $s+r$ & \textbf{Speedup} ($r=20$) & Speedup ($r=50$) & Feasible obj.\ gap (\%)\\
\midrule
2.45  & 2.14  & 42{,}774      & 20{,}017  & \textbf{1.17} & 0.98 & 0.03\\
3.49  & 3.37  & 229{,}175     & 67{,}986  & \textbf{4.84} & 1.95 & 0.02\\
8.64  & 7.77  & 709{,}567     & 91{,}297  & \textbf{1.00} & 0.66 & 0.72\\
11.09 & 9.84  & 903{,}516     & 91{,}780  & \textbf{0.74} & 0.67 & 1.18\\
15.34 & 13.21 & 1{,}323{,}999 & 100{,}185 & \textbf{1.07} & 0.67 & 1.88\\
\bottomrule
\end{tabular}}
\end{table}

Two readings. First, the rank rule is worth a factor of about $2.5$ at the peak --- $4.84
\times$ against $1.95\times$ --- at accuracy identical to two decimal places, so a
substantial part of the attainable speedup in the submitted study was lost to the rank
choice alone. Second, and unlike the grid, the speedup on Chicago Sketch is \emph{not}
monotone in richness: it peaks near $\bar K\approx3.5$ and settles close to unity beyond
$\bar K\approx9$, while the feasible gap grows from $0.03\%$ to $1.88\%$. Equation
\eqref{eq:cost-rule} explains the difference between the two families. Chicago Sketch paths
average about $15$ links, so even at $r=20$ the dense term $2r=40$ exceeds the sparse term;
grid paths at $N=32$ average $35$--$45$ links, and the inequality runs the other way. The
compression is thus best suited to networks with long paths and many parallel routes, and the
enriched-pool regime of a short-path network is its adverse case.

Applying the rank rule to the submitted threshold grid improves accuracy as well as time: at
$r=20$ the conversion $\delta_F$ vanishes at every threshold, the feasible objective gaps are
$0.027\%$, $0.000\%$, $0.020\%$ and $0.028\%$ for $\tau=0$, $0.46$, $1.06$ and $4.54$, and
the link-flow $R^2$ rises to between $0.9997$ and $1.0000$, compared with $0.990$--$0.996$ at
$r=50$. The distribution of demand between major and minor paths at each threshold is
reported in Appendix~\ref{app:pathflow}.

\conote{Chicago Regional and Philadelphia are not yet re-evaluated at the corrected rank. The
submitted Regional pool ($879{,}625$ paths) is no longer on file; the closest available
instance has $2{,}024{,}525$ paths and $30$ links per path, which by \eqref{eq:cost-rule}
should favour $r\approx10$. A C\texttt{++} rank-by-offset experiment on it is running.
Referee~1 will reasonably note that our largest speedups come from Sioux Falls and a synthetic
grid, so carrying at least one large road network at the corrected rank matters more than any
other remaining item.}


\subsection{Reading the Two Representational Choices Together}
\label{sec:tradeoff}

The subspace criterion and the affine offset act on different parts of the same representation,
and both were measured by the same device: an instance solved twice, back to back, with one
choice changed and everything else held fixed. Figure~\ref{fig:paired} shows every such pair as
a displacement rather than as a point, which is the only form in which these comparisons are
controlled. The absolute position of a configuration in the speed--accuracy plane is dominated
by which instance it is; the displacement isolates the choice under study.

\begin{figure}[h]
\centering
\includegraphics[width=\textwidth]{fig/fig_paired.pdf}
\caption{Paired comparisons. Each arrow is one instance solved twice with a single change,
running from the plain object to the one advocated here: (a) plain to flow-weighted subspace,
(b) no offset to offset present. Open markers are the baseline. The horizontal axis is exact;
the vertical axis carries the measurement noise of Section~\ref{sec:setup}, drawn as the error
bar, so vertical displacements shorter than it are not effects.}
\label{fig:paired}
\end{figure}

Comparing the two panels separates the two mechanisms cleanly, and the noise floor is what
makes the separation legible. In panel~(a) the arrows point left in $27$ of $34$ pairs, so the
weighted subspace improves accuracy, and that is measured exactly; vertically they scatter, and
only $2$ of $34$ rise by more than the noise floor, so we claim no speed effect for the
weighting. Panel~(b) behaves differently. Adding the offset improves the gap in $16$ of $20$
grid cells \emph{and} improves the solution time by more than the noise floor in $10$ of $20$.

The offset is therefore a speed mechanism as well as an accuracy one, whereas the subspace
criterion is an accuracy mechanism only. This is what one would expect from
Sections~\ref{sec:weighted-basis} and \ref{sec:offset-study}: weighting rotates the retained
subspace, changing what the representation can express but not how hard the subproblem is,
while the offset moves the starting point into the interior of the feasible set and so changes
the difficulty of the subproblem directly. The two act on different parts of the representation
and neither substitutes for the other.

Figure~\ref{fig:3d} shows the same two comparisons over the full grid design, on the axis that
can be resolved. The weighted sheet lies below the plain one and the offset-free sheet lies
above both, and the separation grows with pool depth, which is where the criterion has the most
to gain.

\begin{figure}[h]
\centering
\includegraphics[width=\textwidth]{fig/fig_3d.pdf}
\caption{Feasible objective gap over the grid design $(N,K)$. (a) plain (orange) against
flow-weighted (blue): the weighted sheet is lower in $15$ of $20$ cells. (b) offset present
(blue) against removed (red): removing the anchor lifts the surface in $16$ of $20$ cells.}
\label{fig:3d}
\end{figure}



\subsection{The Repeated-Scenario Setting, and When Preprocessing is Repaid}
\label{sec:repeated}

The framework is aimed at settings where closely related assignments are solved repeatedly, and
two objections turn on whether that setting is real. The nominal flow is obtained by solving the
full problem, so compressing a single instance is circular; and a path-based method could
simply warm-start from the previous solution and obtain the same benefit without any
compression. Both objections concern the same experiment, so we run it.

On the grid family, a baseline scenario is solved uncompressed. Its solution supplies two
things: the nominal flow $w_0$ from which the basis is built, and the warm start for the
competing method. Each follow-up scenario perturbs either the demand or the link capacities by
a relative amount $\delta$ drawn uniformly in $[-\delta,\delta]$, with the candidate path pool
and the basis held \emph{fixed} at their baseline values --- nothing is recomputed. Each
scenario is then solved three ways: cold and uncompressed, to serve as the accuracy reference;
uncompressed but warm-started from the baseline solution, which is the competing proposal; and
compressed, reusing the baseline basis unchanged.

\begin{table}[h]
\centering
\caption{Repeated-scenario experiment on the grid ($N=32$, $K=64$, $n=66{,}161$ paths, $86.4\%$
variable reduction, matched tolerance). The pool and the basis are built once from the baseline
and never rebuilt. Gaps are measured against a cold uncompressed solve of the same scenario.}
\label{tab:repeated}
{\small
\begin{tabular}{llrrrrrr}
\toprule
& & \multicolumn{1}{c}{cold} & \multicolumn{2}{c}{warm-started uncompressed}
& \multicolumn{3}{c}{compressed, reused basis}\\
\cmidrule(lr){3-3}\cmidrule(lr){4-5}\cmidrule(lr){6-8}
perturbation & $\delta$ & CPU (s) & CPU (s) & gap (\%) & CPU (s) & gap (\%) & link err.\ (\%)\\
\midrule
demand   & 0.05 & 52.70 & 32.77 & 0.0053 & \textbf{12.60} & 0.0000 & 5.63\\
demand   & 0.10 & 71.77 & 54.28 & 0.0083 & \textbf{14.55} & 0.0000 & 6.32\\
demand   & 0.20 & 65.55 & 37.66 & 0.0095 & \textbf{11.46} & 0.0000 & 6.24\\
capacity & 0.05 & 61.64 & 54.16 & 0.0115 & \textbf{13.39} & 0.0000 & 5.25\\
capacity & 0.10 & 67.10 & 57.79 & 0.0207 & \textbf{14.20} & 0.0000 & 5.05\\
capacity & 0.20 & 67.20 & 63.70 & 0.0381 & \textbf{18.55} & 0.0000 & 4.24\\
\midrule
mean     &      & 64.33 & 50.06 & 0.0156 & \textbf{14.12} & 0.0000 & 5.45\\
\bottomrule
\end{tabular}}
\end{table}

Three things follow from Table~\ref{tab:repeated}. First, warm-starting does help the
uncompressed method --- it cuts the mean solve from $64.3$s to $50.1$s --- so the objection is
well founded, but it does not close the gap: the compressed solve averages $14.1$s, a factor of
$3.5$ faster than the warm-started full method and comfortably beyond the noise floor. Second,
the compressed solve is not paying for this in objective accuracy. Its feasible objective gap
is $0.0000\%$ in every scenario, better than the warm-started full solve, whose gap grows with
the perturbation as it inherits a stale starting point. Third, it \emph{is} paying at the link
level: the compressed link-flow error averages $5.45\%$ against roughly $2\%$ for the
warm-started method. A representation built for one scenario reproduces the objective of a
nearby scenario very well while distributing flow somewhat differently across links, which
matters for any application that reads per-link levels of service and is stated here rather
than left for the reader to discover.

\paragraph{Break-even.} The preprocessing is one uncompressed solve of the baseline plus the
factorisation: $65.22 + 0.76 = 65.98$s, of which the basis is only $1\%$ --- the cost is
essentially the nominal flow, not the SVD. With a saving of $50.06 - 14.12 = 35.94$s per
follow-up scenario, the number of related solves needed to repay it is
\begin{equation}
\label{eq:breakeven}
H^\ast=\left\lceil\frac{T_{\text{prep}}}{T_{\text{warm}}-T_{\text{compressed}}}\right\rceil
=\left\lceil\frac{65.98}{35.94}\right\rceil=2 .
\end{equation}
Two follow-up scenarios repay the entire preprocessing cost, and every scenario after that is
a net saving; had the baseline solve been available already --- which it is whenever the
compression is applied to an existing calibrated model --- the basis alone costs $0.76$s and
$H^\ast=1$. This is the quantitative form of the claim that the method is for repeated
assignment, and it is the answer to the objection that compressing a single instance is
circular: for one solve it is, and the framework does not claim otherwise.

\subsection{Scope and Limitations}
\label{sec:scope-limits}


Four limitations bound the claims above, and each is a measurement rather than a caveat.

\emph{The speedup depends on the implementation.} On the richness axis of
Table~\ref{tab:cutoff-summary} the Python implementation peaks at moderate richness and falls
away, whereas the compiled implementation rises monotonically towards unity without crossing
it; the compiled solver's sparse kernels are relatively faster than its dense ones, which
moves the balance in \eqref{eq:cost-rule}. Sioux Falls is favourable in both ($3.6\times$ and
$6.2\times$). Speedup is therefore a property of representation, operator, implementation and
instance geometry jointly, and we report it only within an implementation.

\emph{The dense reconstruction is the binding constraint.} Reducing the number of independent
variables does not reduce the $n-s$ reconstructed inequalities. This is what
\eqref{eq:cost-rule} formalizes, what the falling branch of Table~\ref{tab:cutoff-summary}
measures, and what limits the method on short-path networks with rich pools.

\emph{A grouped representation removes that limit but is not yet faster.} Replacing the
signed subspace by nonnegative groups makes feasibility a bound constraint by construction,
with no reconstruction at all. Measured on the adverse cells of
Table~\ref{tab:cutoff-summary}, it is essentially exact --- feasible gaps of $0.016\%$ and
$0.000\%$, in one case improving on the uncompressed solution --- but in the present
implementation it is not faster. It is the natural representation for the regime where the
signed subspace struggles, and we leave its efficient implementation to future work.

\emph{Timing precision.} Large-instance timings are single runs; repeated identical solves
varied by up to $40\%$ under machine load, so ratios near unity on those instances should not
be read as ordered.


\end{document}
